% Generated from the matching Typst scaffold by
% scripts/migrate_typst_report_to_latex.py.

% ── 5.3  H2: Does What You Forget Matter? ────────────────────────────────
% PURPOSE: Isolate whether content-aware forgetting beats blind volume reduction by
%          pitting Type-Aware Decay directly against Random Prune at the same cap.
% MOVE: Pain → Consequence → Gap (volume control is cheap; does content selection earn
%        its complexity?).
% FIGURES/TABLES: tab:main-results (Type-Aware vs Random subset).
% CITATIONS: anderson1991 (the multiplicative decay rationale behind Type-Aware).
% ─────────────────────────────────────────────────────────────────────────

\section{H2: Does What You Forget Matter?}
\label{sec:h2}

The volume question (\autoref{sec:h1}) settles only \emph{how much} memory an agent
should keep, not \emph{which} items it should keep. H2 asks the sharper question: at a
fixed budget, does selecting \emph{what} to forget by content beat discarding items at
random? The comparison is deliberately clean. Type-Aware Decay and Random Prune both
evict to the same capacity, both leave retrieval scoring untouched (Invariant~5), and both
hold the embedding payload constant (Invariant~4); the only thing that varies between them
is the eviction rule. Random Prune drops a uniformly chosen item; Type-Aware Decay scores
each item with the multiplicative, Anderson--Schooler-style rule of Invariant~8
\citep{anderson1991} --- a per-type age decay tempered by a retrieval-frequency term ---
and evicts the lowest-scoring item. Any residual gap in CL-F1 is therefore attributable to
content selection, not to footprint. The question is whether that machinery earns its
complexity over a coin flip.

It does not. As reported in \autoref{tab:h2-content-vs-random}, content-aware pruning does
not beat random pruning on the trustworthy 144-run matrix. Type-Aware Decay attains the
\emph{lowest} CL-F1 of all six policies (0.427), below Random Prune (0.447), below CLS
Consolidation (0.449), and --- the striking part --- below even the No-Memory baseline
(0.429). Recall that CL-F1 here is the resolved-rate proxy of \autoref{sec:metrics}, not a
forgetting-curve measure; the policies are compared on the proxy throughout this chapter.
The signed-rank contrast of Type-Aware Decay against full memory (the reference used for
the pre-registered family of five comparisons in \autoref{sec:h1}) is not significant
($W=10.0$, $p=0.297$, Holm-adjusted $p=1.000$; median $\Delta=-0.0024$; rank-biserial
$r_{\mathrm{rb}}=-0.444$, a medium effect in the \emph{negative} direction; $95\%$ BCa CI
$[-0.063,\,0.016]$, which crosses zero). The confidence interval is bounded below at a
modest magnitude, so we read this as no detectable content-over-volume benefit rather than
a strong harm; we do not run a TOST equivalence verdict here (it was pre-registered but not
computed). The plain reading is that at a fixed budget, \emph{which} item you forget does
not move the resolved-rate proxy --- volume reduction alone is the active ingredient, and a
coin flip selects as well as the theoretically motivated rule.

\begin{table}[htbp]
\centering
\caption{H2 content-vs-random comparison. CL-F1 is the resolved-rate proxy averaged over
the eight sequences; both pruners hold footprint at the same cap and share identical
retrieval scoring, so the contrast isolates content selection. Type-Aware Decay does not
beat Random Prune, and falls below the No-Memory baseline. Wilcoxon statistics are reported
against the full-memory reference (Invariant~11), with rank-biserial $r_{\mathrm{rb}}$ and
$95\%$ BCa CI.}
\label{tab:h2-content-vs-random}
\begin{tabular}{lcccc}
\toprule
Policy & CL-F1 (resolve) & Footprint (tok/task) & $r_{\mathrm{rb}}$ vs full & Holm $p$ \\
\midrule
CLS Consolidation & 0.449 & 704 & $-0.071$ & 1.000 \\
Random Prune      & 0.447 & 601 & $-0.056$ & 1.000 \\
Full Memory       & 0.445 & 1{,}247 & --- & --- \\
Recency Prune     & 0.437 & 604 & $-0.111$ & 1.000 \\
No Memory         & 0.429 & 0 & $-0.286$ & 1.000 \\
Type-Aware Decay  & \textbf{0.427} & 598 & $-0.444$ & 1.000 \\
\bottomrule
\end{tabular}
\end{table}

\subsection{The striking negative: Type-Aware Decay freezes its own memory}
\label{sec:h2-tad-mechanism}

A policy designed to forget intelligently performing \emph{worse} than keeping no memory at
all demands a mechanism, not just a number. A content audit of the Type-Aware Decay runs
locates it in the eviction score itself. The Invariant~8 rule weighs a per-type age decay
against a retrieval-frequency term of the form $(1+\text{use\_count})^{0.5}$. In practice
the retrieval-frequency term dominates the age-decay term: once an early item has been
retrieved a few times, its score is buoyed faster than age erodes it, so it is never the
lowest-scoring candidate and is never evicted. Newly written items, with a use count of
zero, are the cheapest to discard and are evicted almost immediately. The net effect is
that the store \emph{freezes} at roughly its first ten items --- the early high-use entries
become effectively permanent while every later memory is created and destroyed within a
task or two. The agent is left carrying a small, stale core that rarely matches the task at
hand, which is plausibly why Type-Aware Decay underperforms even the No-Memory baseline:
the No-Memory agent at least pays no retrieval-noise tax, whereas the frozen store
sometimes injects irrelevant early context. This is a genuine negative result, not a tuning
artifact: the decay parameters were frozen at design time per the pre-registration and were
never re-fit from the data (Invariant~8), so the failure is a property of the formula as
specified rather than of a poorly chosen knob.

Taken together, H2 finds no support for the hypothesis that content-aware forgetting beats
blind volume reduction. At a matched budget, random eviction is as good as the
Anderson--Schooler-style rule, and the one policy that selects most aggressively by content
is the worst performer in the matrix. The parsimonious reading --- consistent with the
powered policy-null of \autoref{sec:h1} --- is that on this resolved-rate proxy the active
ingredient is footprint, not the cleverness of the selection rule.
